\scanpage{7}
\mspage{\textcircled{1}}
% ll. 1–5
$X\in\operatorname{Ob}\mathrm{Cat}$ tel que tout objet \unsure{0-connexe|0-conn.} de $X^\wedge$ soit 1-connexe

$F\in\operatorname{Ob} X^\wedge$, $U=\operatorname{Im}(F\longrightarrow e_{X^\wedge})$. Donc $F\in\operatorname{Ob}X^\wedge_{/U}\unsure{\simeq|\cong}(\underbrace{X_{/U}}_Y)^\wedge=\deleted{\unsure{X|x}}Y^\wedge$, et en \deleted{\unsure{tant|tel}} qu'objet de $Y^\wedge$, $F\longrightarrow e_{Y^\wedge}$ épi.

% ll. 6–8
\deleted{\unsure{$I$|Soit}} Soit $\widetilde\Delta={\rm Cat.}$ des ens. finis de la forme $\widetilde\Delta_n=[0,n]\cap\mathbb Z$ (équiv. à la cat. des ens finis non vides)
% ll. 9–11
\[
\begin{aligned}
 I&\longmapsto\pi_0(F^I)\ :\ \widetilde\Delta^{\mathrm o}\longrightarrow(\mathrm{Ens}),\quad\text{soit}\\
 \pi_F&=\text{\uline{$\pi\in\widetilde\Delta^\wedge$}}\ :\qquad\pi(I)=\pi_0(F^I)\qquad\text{i.e.}\\
 &&\pi(\widetilde\Delta_{\unsure{n|m}})=\pi_0(F^{n+1}\deleted{_{\unsure{*|1}}})
\end{aligned}
\]
% ll. 12–14
Je vais définir un foncteur
\[
 \boxed{\unsure{\varphi|\Phi}:(\widetilde\Delta_{/\pi})^{\mathrm o}\longrightarrow Y^\wedge=(X_{/U})^\wedge\approx X^\wedge_{/U}}
\]
% ll. 15–18
Un objet de $\widetilde\Delta_{/\pi}$ est un couple $(\widetilde\Delta_n,\alpha)$, avec $\alpha\in\pi(\widetilde\Delta_n)=\pi_0(F^{n+1})$, on pose
\[
 \varphi(\widetilde\Delta_n,\alpha)=\underbrace{\alpha}_{\substack{\text{composante}\\\text{conn. de }F^{n+1}}}\subset\deleted{\unsure{e|\alpha}}F^{n+1}
\]
\margin{Comme $F$, les $F^I$, sont ds $X^\wedge_{/U}\overset{\approx}{\longrightarrow}Y^\wedge$ de même $\alpha$.}
% ll. 19–22
Une flèche de $\widetilde\Delta_{/\pi}$ \deleted{est}, de $(\widetilde\Delta_n,\alpha)$ vers $(\widetilde\Delta_m,\beta)$, est formé d'une flèche $f:\widetilde\Delta_n\longrightarrow\widetilde\Delta_m$ telle que $\underset{\substack{\rotatebox{90}{$=$}\ {\rm def}\\\pi(\rho)(\beta)}}{{\rho}_\pi^*(\beta)}=\alpha$. Or
% archive sheet number: 6
